\documentclass[12pt,a4paper]{article}


\usepackage[utf8]{inputenc}
\usepackage{xeCJK}
\usepackage{geometry}
\usepackage{tikz}
\usepackage{listings}
\usepackage{xcolor}
\usepackage{hyperref}
\usepackage{enumitem}
\usepackage{booktabs}
\usepackage{longtable}

\geometry{margin=2.5cm}
\setCJKmainfont{Noto Sans CJK TC}

% 設定程式碼樣式
\lstset{
    basicstyle=\ttfamily\small,
    backgroundcolor=\color{gray!10},
    frame=single,
    breaklines=true,
    showstringspaces=false,
    numbers=left,
    numberstyle=\tiny\color{gray}
}

% 這是一份latex 文件的 preamble
% 請AI閱讀後，將附上的PDF整理，以latex code整理，儘量用longtable
% 請注意我的名字: 謝慕揚, MD, PhD, FESC
% 表格請注意換行，不該在cell內使用 \\
% 表格請不要有垂直分隔線 |



% Including essential packages for Chinese typesetting
\usepackage{xeCJK}
\usepackage{fontspec}
% Configuring Chinese font with Noto Serif CJK TC, fallback to SimSun
\IfFontExistsTF{Noto Serif CJK TC}{
  \setCJKmainfont{Noto Serif CJK TC}
  \setCJKsansfont{Noto Serif CJK TC}
  \setCJKmonofont{Noto Serif CJK TC}
}{\setCJKmainfont{SimSun}
  \setCJKsansfont{SimSun}
  \setCJKmonofont{SimSun}
}

% Including other necessary packages
\usepackage{geometry}
\usepackage{titlesec}
\usepackage{enumitem}
\usepackage{xcolor}
\usepackage{colortbl}
\usepackage{tcolorbox}
\usepackage{booktabs}
\usepackage{longtable}
\usepackage{array}
\usepackage{graphicx}
\usepackage{tikz}
\usepackage{fancyhdr}
\usepackage{multicol}
\usepackage{datetime2}
\usepackage{hyperref}
\usepackage{mdframed}
\usepackage{amsmath}

\hypersetup{
    colorlinks=true,
    linkcolor=emergency_blue,
    citecolor=emergency_blue,
    urlcolor=emergency_blue,
    pdfborder={0 0 0}
}

% Defining page geometry
\geometry{
  top=2.5cm,
  bottom=2.5cm,
  left=2cm,
  right=2cm,
  headheight=25pt
}

% Adjusting topmargin to compensate for increased headheight
\addtolength{\topmargin}{-10pt}

% Defining colors
\definecolor{emergency_red}{RGB}{186,24,27}
\definecolor{emergency_blue}{RGB}{0,114,188}
\definecolor{emergency_gray}{RGB}{120,120,120}
\definecolor{highlight_yellow}{RGB}{255,250,205}

% Setting title formats
\titleformat{\section}[block]{\Large\bfseries\color{emergency_red}}{\thesection}{10pt}{}
\titleformat{\subsection}[block]{\large\bfseries\color{emergency_blue}}{\thesubsection}{8pt}{}
\titleformat{\subsubsection}[block]{\normalsize\bfseries\color{emergency_gray}}{\thesubsubsection}{6pt}{}

% Defining custom environment for important notes
\newmdenv[
    linecolor=emergency_red,
    backgroundcolor=highlight_yellow,
    linewidth=2pt,
    topline=true,
    bottomline=true,
    leftline=true,
    rightline=true,
    innertopmargin=10pt,
    innerbottommargin=10pt,
    innerleftmargin=10pt,
    innerrightmargin=10pt
]{importantbox}

% Defining note command with tcolorbox
\newcommand{\note}[1]{
    \par
    \noindent
    \begin{tcolorbox}[
        colback=emergency_blue!5!white,
        colframe=emergency_blue,
        title=\textbf{重點提醒},
        fonttitle=\bfseries,
        boxrule=0.5pt,
        arc=3pt,
        left=6pt,
        right=6pt,
        top=6pt,
        bottom=6pt
    ]
    #1
    \end{tcolorbox}
}

% % Setting up header and footer
% \pagestyle{fancy}
% \fancyhf{}
% \fancyhead[L]{\textcolor{emergency_red}{心血管中心教學文件}}
% \fancyhead[R]{\textcolor{emergency_blue}{急性心肌梗塞早期診斷與心電圖教學}}
% \fancyfoot[C]{\textcolor{emergency_gray}{\thepage}}
% \renewcommand{\headrulewidth}{1pt}
% \renewcommand{\headrule}{\hbox to\headwidth{\color{emergency_red}\leaders\hrule height \headrulewidth\hfill}}

% Defining custom date format for ISO 8601
\DTMnewdatestyle{iso}{
  \renewcommand{\DTMdisplaydate}[4]{\number##1-\number##2-\number##3}
  \renewcommand{\DTMDisplaydate}[4]{\number##1-\number##2-\number##3}
}
\DTMsetdatestyle{iso}
\newcommand{\mydate}{\DTMtoday}


\title{\textbf{YouTube視頻轉文字稿生成講義系統}}
\author{自動化教學內容生產流程}
\date{\today}

\begin{document}

\maketitle
\tableofcontents
\newpage

\section{系統概述}

本系統提供多種方法將YouTube視頻轉換為高品質的文字稿，並透過AI處理生成結構化的教學講義。適用於線上課程、研討會、講座等教育內容的快速整理。

\section{方法比較表}

\begin{longtable}{|p{3cm}|p{3cm}|p{3cm}|p{3cm}|p{2cm}|}
\hline
\textbf{方法} & \textbf{優點} & \textbf{缺點} & \textbf{適用場景} & \textbf{成本} \\
\hline
YouTube自動字幕 & 免費、快速 & 準確度有限 & 英文內容較佳 & 免費 \\
\hline
Whisper AI & 高準確度、多語言 & 需技術設定 & 所有語言 & 免費 \\
\hline
商業工具 & 操作簡單、功能完整 & 需要付費 & 專業使用 & 付費 \\
\hline
手動+AI輔助 & 最高品質 & 時間較長 & 重要內容 & 中等 \\
\hline
\end{longtable}

\section{完整實施流程}

\begin{center}
\begin{tikzpicture}[node distance=2.5cm, auto]
    % 定義節點樣式
    \tikzstyle{process} = [rectangle, rounded corners, minimum width=3.5cm, minimum height=1cm, text centered, draw=black, fill=blue!20]
    \tikzstyle{decision} = [diamond, minimum width=3cm, minimum height=1cm, text centered, draw=black, fill=orange!20]
    \tikzstyle{io} = [trapezium, trapezium left angle=70, trapezium right angle=110, minimum width=3cm, minimum height=1cm, text centered, draw=black, fill=green!20]
    \tikzstyle{arrow} = [thick,->,>=stealth]

    % 流程節點
    \node (start) [io] {YouTube視頻URL};
    \node (extract) [process, below of=start] {音頻提取\\(yt-dlp)};
    \node (transcribe) [process, below of=extract] {語音轉文字\\(Whisper AI)};
    \node (clean) [process, below of=transcribe] {文字清理\\格式化};
    \node (ai) [process, below of=clean] {AI內容整理\\結構化};
    \node (latex) [process, below of=ai] {LaTeX講義\\生成};
    \node (output) [io, below of=latex] {PDF講義\\完成};

    % 流程箭頭
    \draw [arrow] (start) -- (extract);
    \draw [arrow] (extract) -- (transcribe);
    \draw [arrow] (transcribe) -- (clean);
    \draw [arrow] (clean) -- (ai);
    \draw [arrow] (ai) -- (latex);
    \draw [arrow] (latex) -- (output);
\end{tikzpicture}
\end{center}

\section{方法一：使用Whisper AI（推薦）}

\subsection{環境設置}

\begin{lstlisting}[language=bash, caption=安裝必要工具]
# 安裝Python套件
pip install openai-whisper
pip install yt-dlp
pip install pydub

# 安裝FFmpeg（Windows用戶）
# 下載並安裝 https://ffmpeg.org/download.html

# 安裝FFmpeg（Mac用戶）
brew install ffmpeg

# 安裝FFmpeg（Linux用戶）
sudo apt update
sudo apt install ffmpeg
\end{lstlisting}

\subsection{自動化腳本}

\begin{lstlisting}[language=python, caption=YouTube轉文字完整腳本]
import whisper
import yt_dlp
import os
import re
from datetime import datetime

class YouTubeToTranscript:
    def __init__(self):
        # 載入Whisper模型（medium模型平衡速度和準確度）
        self.model = whisper.load_model("medium")
    
    def download_audio(self, youtube_url, output_path="temp_audio"):
        """下載YouTube音頻"""
        ydl_opts = {
            'format': 'bestaudio/best',
            'outtmpl': f'{output_path}.%(ext)s',
            'postprocessors': [{
                'key': 'FFmpegExtractAudio',
                'preferredcodec': 'wav',
                'preferredquality': '192',
            }],
        }
        
        with yt_dlp.YoutubeDL(ydl_opts) as ydl:
            # 獲取視頻資訊
            info = ydl.extract_info(youtube_url, download=False)
            title = info.get('title', 'Unknown')
            duration = info.get('duration', 0)
            
            print(f"正在下載：{title}")
            print(f"時長：{duration//3600:02d}:{(duration%3600)//60:02d}:{duration%60:02d}")
            
            # 下載音頻
            ydl.download([youtube_url])
        
        return f"{output_path}.wav", title
    
    def transcribe_audio(self, audio_path, language="zh"):
        """使用Whisper轉錄音頻"""
        print("正在進行語音識別...")
        result = self.model.transcribe(
            audio_path, 
            language=language,
            task="transcribe",
            word_timestamps=True
        )
        return result
    
    def format_transcript(self, result):
        """格式化轉錄結果"""
        transcript = ""
        
        for segment in result["segments"]:
            start_time = self.seconds_to_timestamp(segment["start"])
            text = segment["text"].strip()
            transcript += f"[{start_time}] {text}\n\n"
        
        return transcript
    
    def seconds_to_timestamp(self, seconds):
        """將秒數轉換為時間戳"""
        hours = int(seconds // 3600)
        minutes = int((seconds % 3600) // 60)
        seconds = int(seconds % 60)
        return f"{hours:02d}:{minutes:02d}:{seconds:02d}"
    
    def clean_transcript(self, transcript):
        """清理和優化文字稿"""
        # 移除重複的語氣詞
        cleaned = re.sub(r'\b(呃|嗯|那個|這個|就是說)\b', '', transcript)
        # 合併短句
        cleaned = re.sub(r'\n\n(?=\[[0-9]{2}:[0-9]{2}:[0-9]{2}\])', '\n', cleaned)
        # 整理標點符號
        cleaned = re.sub(r'([。！？])\s*', r'\1\n\n', cleaned)
        
        return cleaned
    
    def process_youtube_video(self, youtube_url, language="zh", output_file=None):
        """處理YouTube視頻的主要流程"""
        try:
            # 1. 下載音頻
            audio_path, title = self.download_audio(youtube_url)
            
            # 2. 語音轉文字
            result = self.transcribe_audio(audio_path, language)
            
            # 3. 格式化文字稿
            transcript = self.format_transcript(result)
            
            # 4. 清理文字稿
            cleaned_transcript = self.clean_transcript(transcript)
            
            # 5. 保存結果
            if output_file is None:
                output_file = f"{re.sub(r'[<>:\"/\\|?*]', '_', title)}_transcript.txt"
            
            with open(output_file, 'w', encoding='utf-8') as f:
                f.write(f"視頻標題：{title}\n")
                f.write(f"轉錄時間：{datetime.now().strftime('%Y-%m-%d %H:%M:%S')}\n")
                f.write("="*50 + "\n\n")
                f.write(cleaned_transcript)
            
            # 6. 清理臨時檔案
            if os.path.exists(audio_path):
                os.remove(audio_path)
            
            print(f"轉錄完成！文字稿已保存至：{output_file}")
            return output_file, title, cleaned_transcript
            
        except Exception as e:
            print(f"處理過程中發生錯誤：{str(e)}")
            return None, None, None

# 使用範例
if __name__ == "__main__":
    converter = YouTubeToTranscript()
    
    # 輸入YouTube URL
    url = input("請輸入YouTube視頻URL：")
    language = input("請輸入語言代碼（zh=中文, en=英文, auto=自動檢測）：") or "auto"
    
    # 處理視頻
    output_file, title, transcript = converter.process_youtube_video(url, language)
    
    if output_file:
        print(f"\n成功生成文字稿：{output_file}")
    else:
        print("轉錄失敗，請檢查URL或網路連接")
\end{lstlisting}

\section{方法二：使用現有字幕}

\subsection{提取YouTube字幕}

\begin{lstlisting}[language=python, caption=提取YouTube自動字幕]
import yt_dlp
import json

def extract_youtube_subtitles(youtube_url, lang='zh-TW'):
    """提取YouTube自動生成的字幕"""
    ydl_opts = {
        'writesubtitles': True,
        'writeautomaticsub': True,
        'subtitleslangs': [lang, 'zh', 'en'],
        'skip_download': True,
        'subtitlesformat': 'vtt',
    }
    
    with yt_dlp.YoutubeDL(ydl_opts) as ydl:
        try:
            info = ydl.extract_info(youtube_url, download=False)
            title = info.get('title', 'Unknown')
            
            # 檢查是否有字幕
            subtitles = info.get('subtitles', {})
            auto_subtitles = info.get('automatic_captions', {})
            
            available_subs = {**subtitles, **auto_subtitles}
            
            if available_subs:
                print(f"找到字幕語言：{list(available_subs.keys())}")
                
                # 下載字幕
                ydl.download([youtube_url])
                return title, True
            else:
                print("未找到可用字幕")
                return title, False
                
        except Exception as e:
            print(f"提取字幕時發生錯誤：{str(e)}")
            return None, False

# 使用範例
url = "YOUR_YOUTUBE_URL"
title, has_subtitles = extract_youtube_subtitles(url)
\end{lstlisting}

\section{AI處理生成講義}

\subsection{ChatGPT/Claude處理提示詞}

\begin{lstlisting}[caption=AI整理提示詞範例]
請將以下YouTube視頻文字稿整理成結構化的教學講義，要求：

1. **內容結構化**：
   - 提取主要主題和子主題
   - 按邏輯順序重新組織內容
   - 生成清楚的章節標題

2. **重點摘要**：
   - 識別關鍵概念和定義
   - 提取重要的實例和案例
   - 整理數據和統計資訊

3. **LaTeX格式輸出**：
   - 使用適當的章節結構
   - 包含數學公式（如有）
   - 添加圖表說明（如提到）
   - 生成參考文獻格式

4. **學習輔助**：
   - 生成學習目標
   - 添加重點提示框
   - 創建思考問題
   - 製作摘要表格

文字稿內容：
[在此貼上轉錄的文字稿]

請輸出完整的LaTeX代碼，可直接在Overleaf中編譯。
\end{lstlisting}

\subsection{Python自動化AI處理}

\begin{lstlisting}[language=python, caption=自動AI處理腳本]
import openai
import anthropic
from typing import Optional

class TranscriptToNotes:
    def __init__(self, api_key: str, service: str = "openai"):
        self.service = service
        if service == "openai":
            openai.api_key = api_key
        elif service == "anthropic":
            self.client = anthropic.Anthropic(api_key=api_key)
    
    def process_with_ai(self, transcript: str, topic: str = "") -> str:
        """使用AI處理文字稿生成講義"""
        
        prompt = f"""
請將以下關於「{topic}」的視頻文字稿整理成專業的教學講義。

要求：
1. 生成LaTeX格式的完整文件
2. 包含：標題、摘要、章節、重點、結論
3. 提取關鍵概念並加以解釋
4. 整理實例和案例研究
5. 生成學習檢核清單

文字稿：
{transcript}

請輸出可直接編譯的LaTeX代碼：
        """
        
        if self.service == "openai":
            response = openai.ChatCompletion.create(
                model="gpt-4",
                messages=[
                    {"role": "system", "content": "你是一位專業的教學設計師，擅長將原始內容轉換為結構化的教學材料。"},
                    {"role": "user", "content": prompt}
                ],
                max_tokens=4000,
                temperature=0.7
            )
            return response.choices[0].message.content
            
        elif self.service == "anthropic":
            response = self.client.messages.create(
                model="claude-3-sonnet-20240229",
                max_tokens=4000,
                messages=[{"role": "user", "content": prompt}]
            )
            return response.content[0].text
    
    def generate_notes(self, transcript_file: str, topic: str = "", output_file: str = None):
        """從文字稿檔案生成講義"""
        # 讀取文字稿
        with open(transcript_file, 'r', encoding='utf-8') as f:
            transcript = f.read()
        
        # AI處理
        latex_content = self.process_with_ai(transcript, topic)
        
        # 保存結果
        if output_file is None:
            output_file = transcript_file.replace('.txt', '_notes.tex')
        
        with open(output_file, 'w', encoding='utf-8') as f:
            f.write(latex_content)
        
        print(f"講義已生成：{output_file}")
        return output_file

# 使用範例
processor = TranscriptToNotes("YOUR_API_KEY", "openai")
notes_file = processor.generate_notes("transcript.txt", "機器學習基礎")
\end{lstlisting}

\section{LaTeX講義模板}

\begin{lstlisting}[language=TeX, caption=YouTube講義LaTeX模板]
\documentclass[12pt,a4paper]{article}
\usepackage[utf8]{inputenc}
\usepackage{xeCJK}
\usepackage{geometry}
\usepackage{amsmath}
\usepackage{graphicx}
\usepackage{hyperref}
\usepackage{tcolorbox}
\usepackage{enumitem}
\usepackage{booktabs}

\geometry{margin=2.5cm}
\setCJKmainfont{Noto Sans CJK TC}

% 定義特殊框架
\newtcolorbox{keypoint}{
    colback=blue!5!white,
    colframe=blue!75!black,
    title=重點提示
}

\newtcolorbox{example}{
    colback=green!5!white,
    colframe=green!75!black,
    title=實例說明
}

\title{[視頻標題] - 學習講義}
\author{基於YouTube視頻自動生成}
\date{\today}

\begin{document}

\maketitle

\section*{課程資訊}
\begin{itemize}
    \item \textbf{原始視頻}：[YouTube URL]
    \item \textbf{時長}：[視頻時長]
    \item \textbf{生成時間}：\today
\end{itemize}

\tableofcontents
\newpage

\section{學習目標}
完成本講義學習後，您將能夠：
\begin{enumerate}
    \item [AI生成的學習目標1]
    \item [AI生成的學習目標2]
    \item [AI生成的學習目標3]
\end{enumerate}

\section{內容摘要}
[AI生成的內容摘要]

\section{主要內容}

\subsection{[主題1]}
[AI整理的內容]

\begin{keypoint}
[重點提示內容]
\end{keypoint}

\begin{example}
[實例說明內容]
\end{example}

\subsection{[主題2]}
[AI整理的內容]

\section{關鍵概念表}
\begin{table}[h]
\centering
\begin{tabular}{@{}ll@{}}
\toprule
\textbf{概念} & \textbf{定義} \\
\midrule
[概念1] & [定義1] \\
[概念2] & [定義2] \\
\bottomrule
\end{tabular}
\caption{重要概念整理}
\end{table}

\section{學習檢核}
\begin{enumerate}
    \item [檢核問題1]
    \item [檢核問題2]
    \item [檢核問題3]
\end{enumerate}

\section{延伸資源}
\begin{itemize}
    \item 原始視頻：\href{[YouTube URL]}{[視頻標題]}
    \item 相關閱讀：[AI建議的參考資料]
\end{itemize}

\end{document}
\end{lstlisting}

\section{商業工具推薦}

\subsection{線上服務}
\begin{itemize}
    \item \textbf{Otter.ai}：專業轉錄服務，支援即時轉錄
    \item \textbf{Rev.com}：高品質人工+AI轉錄
    \item \textbf{Trint}：多語言轉錄和編輯平台
    \item \textbf{Sonix}：自動轉錄和翻譯服務
\end{itemize}

\subsection{桌面軟體}
\begin{itemize}
    \item \textbf{Adobe Premiere Pro}：內建語音轉文字功能
    \item \textbf{DaVinci Resolve}：免費影音編輯軟體，含轉錄功能
    \item \textbf{Descript}：音頻編輯和轉錄一體化工具
\end{itemize}

\section{品質優化建議}

\subsection{轉錄品質提升}
\begin{enumerate}
    \item \textbf{音頻品質}：選擇高品質音頻（192kbps以上）
    \item \textbf{語速適中}：避免過快或過慢的語速
    \item \textbf{清楚發音}：確保講者發音清晰
    \item \textbf{背景降噪}：使用音頻處理軟體降低背景雜音
\end{enumerate}

\subsection{AI處理優化}
\begin{enumerate}
    \item \textbf{分段處理}：長視頻分段處理以提高準確度
    \item \textbf{上下文提供}：給AI提供課程背景和領域資訊
    \item \textbf{術語詞典}：建立專業術語詞典供AI參考
    \item \textbf{人工校對}：AI處理後進行人工校對和修正
\end{enumerate}

\section{完整自動化腳本}

\begin{lstlisting}[language=python, caption=YouTube轉講義完整自動化]
#!/usr/bin/env python3
"""
YouTube視頻轉講義完整自動化系統
"""

import argparse
import os
from youtube_to_transcript import YouTubeToTranscript
from transcript_to_notes import TranscriptToNotes

def main():
    parser = argparse.ArgumentParser(description='YouTube視頻轉講義')
    parser.add_argument('url', help='YouTube視頻URL')
    parser.add_argument('--topic', help='課程主題', default='')
    parser.add_argument('--language', help='語言代碼', default='zh')
    parser.add_argument('--api-key', help='AI API密鑰', required=True)
    parser.add_argument('--service', help='AI服務', choices=['openai', 'anthropic'], default='openai')
    
    args = parser.parse_args()
    
    print("=== YouTube視頻轉講義系統 ===")
    print(f"視頻URL: {args.url}")
    print(f"主題: {args.topic}")
    print(f"語言: {args.language}")
    
    # 步驟1：轉錄視頻
    print("\n[1/3] 正在轉錄視頻...")
    converter = YouTubeToTranscript()
    transcript_file, title, transcript = converter.process_youtube_video(
        args.url, args.language
    )
    
    if not transcript_file:
        print("轉錄失敗，程序結束")
        return
    
    # 步驟2：AI處理生成講義
    print("\n[2/3] 正在生成講義...")
    processor = TranscriptToNotes(args.api_key, args.service)
    topic = args.topic or title
    notes_file = processor.generate_notes(transcript_file, topic)
    
    # 步驟3：完成
    print(f"\n[3/3] 完成！")
    print(f"文字稿: {transcript_file}")
    print(f"講義: {notes_file}")
    print(f"可在Overleaf中編譯LaTeX文件")

if __name__ == "__main__":
    main()
\end{lstlisting}

\section{結論}

透過本系統，您可以：

\begin{itemize}
    \item \textbf{快速轉錄}：10分鐘視頻約需2-3分鐘完成轉錄
    \item \textbf{高品質輸出}：AI整理生成結構化講義
    \item \textbf{多語言支援}：支援中文、英文等多種語言
    \item \textbf{自動化流程}：一鍵完成從視頻到講義的轉換
    \item \textbf{可定制化}：可根據需求調整模板和處理邏輯
\end{itemize}

建議根據您的具體需求選擇合適的方法，並逐步建立自己的自動化工作流程。

\end{document}